\documentclass[12pt]{article}
\usepackage[utf8]{inputenc}
\usepackage[T1]{fontenc}
\usepackage{hyperref}
\title{Alcozauca Mixtec Language Resources for K-12 Curriculum Development}
\author{Compiled by Claude Sonnet 5.5 (claude-sonnet-5-5) for Andr\'as Kornai}
\date{}
\begin{document}
\maketitle

\section{Introduction}
Alcozauca Mixtec (ISO 639-3 xta; Glottocode alco1235) is an Otomanguean language of the Alcozauca municipality in Guerrero, Mexico. The reference description gives about 10,000 speakers and says it is not known to be taught in schools. The searches for this report found very little. The only dedicated reference work is a small dictionary of the closely related Xochapa variety, and the other items are audio and scripture lists and aggregator pages. No school curriculum materials turned up.

\section{Primary Educational Resources}
No primers, textbooks or school programmes for Alcozauca Mixtec were found in the search results.

\section{Linguistic and Reference Materials}
\begin{itemize}
\item \emph{Diccionario b\'asico del mixteco de Xochapa, Guerrero} (Johnson, Gonz\'alez de Guzm\'an and Stark, 2nd ed., 2006). This is a basic Mixtec--Spanish dictionary of about 500 headwords, of the Xochapa variety spoken in Alcozauca municipality. Bibliographic record: \url{https://glottolog.org/resource/reference/id/50994}
\item ASJP word-list record for Mixtec Alcozauca, which lists an earlier 1999 SIL edition of the dictionary: \url{https://asjp.clld.org/languages/MIXTEC_ALCOZAUCA}
\item Glosbe Alcozauca Mixtec--English translations. These are aggregated from various sources, are not a curated dictionary and should be used with caution: \url{https://glosbe.com/xta/en}
\end{itemize}

\section{Religious and Community Resources}
\begin{itemize}
\item Global Recordings Network lists audio recordings in Mixtec, Alcozauca, mainly Bible lessons and evangelistic programmes, and also lists scripture resources for Mixtec, Xochapa: \url{https://globalrecordings.net/en/language/xta}
\end{itemize}

\section{Digital and Audio Resources}
The Global Recordings Network page above is the only audio source found. It is religious in content, so it would need vetting before classroom use.

\section{Curriculum Development Context}
Materials for this variety are very sparse. A curriculum would have to rely on new work with community speakers, using the Xochapa dictionary as a lexical starting point. A discussion thread on a Mixtec forum suggested that Alacatlatzala Mixtec has somewhat more material. That is an opinion from one commenter and has not been checked here.
\end{document}